\documentclass[12pt, letterpaper]{article}

% ── Packages ──────────────────────────────────────────────────────────────────
\usepackage[margin=1in]{geometry}
\usepackage[T1]{fontenc}
\usepackage[utf8]{inputenc}
\usepackage{lmodern}
\usepackage{microtype}
\usepackage{booktabs}
\usepackage{array}
\usepackage{longtable}
\usepackage{tabularx}
\usepackage{xcolor}
\usepackage{listings}
\usepackage{enumitem}
\usepackage{titlesec}
\usepackage{fancyhdr}
\usepackage[most,breakable]{tcolorbox}
\usepackage{hyperref}
\usepackage{parskip}
\usepackage{amsmath}
\usepackage{amssymb}
\usepackage{multicol}
\setlength{\headheight}{14pt}

% ── Colors ────────────────────────────────────────────────────────────────────
\definecolor{warnorange}{RGB}{255,140,0}
\definecolor{noteblue}{RGB}{30,100,200}
\definecolor{codegray}{RGB}{245,245,245}
\definecolor{darkgray}{RGB}{60,60,60}
\definecolor{sectionrule}{RGB}{30,100,200}

% ── Hyperref ──────────────────────────────────────────────────────────────────
\hypersetup{
    colorlinks=true,
    linkcolor=noteblue,
    urlcolor=noteblue,
    pdftitle={Wireless Drive-By-Wire Design Document},
    pdfauthor={Capstone Spring 2026}
}

% ── Code listings ─────────────────────────────────────────────────────────────
\lstdefinestyle{terminal}{
    backgroundcolor=\color{codegray},
    basicstyle=\ttfamily\small,
    breaklines=true,
    frame=single,
    rulecolor=\color{gray},
    framesep=6pt,
    xleftmargin=8pt, xrightmargin=8pt,
    aboveskip=8pt, belowskip=8pt
}
\lstdefinestyle{cpp}{
    backgroundcolor=\color{codegray},
    basicstyle=\ttfamily\small,
    keywordstyle=\color{blue}\bfseries,
    stringstyle=\color{orange!80!black},
    commentstyle=\color{green!50!black}\itshape,
    language=C++,
    breaklines=true,
    frame=single,
    rulecolor=\color{gray},
    framesep=6pt,
    xleftmargin=8pt, xrightmargin=8pt,
    aboveskip=8pt, belowskip=8pt
}
\lstdefinestyle{python}{
    backgroundcolor=\color{codegray},
    basicstyle=\ttfamily\small,
    keywordstyle=\color{blue}\bfseries,
    stringstyle=\color{orange!80!black},
    commentstyle=\color{green!50!black}\itshape,
    language=Python,
    breaklines=true,
    frame=single,
    rulecolor=\color{gray},
    framesep=6pt,
    xleftmargin=8pt, xrightmargin=8pt,
    aboveskip=8pt, belowskip=8pt
}

% ── Callout boxes ─────────────────────────────────────────────────────────────
\newtcolorbox{warningbox}{
    colback=warnorange!12, colframe=warnorange,
    fonttitle=\bfseries, title={$\bigtriangleup$\ WARNING}, breakable
}
\newtcolorbox{notebox}{
    colback=noteblue!8, colframe=noteblue,
    fonttitle=\bfseries, title={Note}, breakable
}
\newtcolorbox{designbox}{
    colback=gray!8, colframe=darkgray,
    fonttitle=\bfseries, title={Design Decision}, breakable
}

% ── Header / Footer ───────────────────────────────────────────────────────────
\pagestyle{fancy}
\fancyhf{}
\fancyhead[L]{\small\textcolor{darkgray}{Wireless Drive-By-Wire}}
\fancyhead[R]{\small\textcolor{darkgray}{Design Document}}
\fancyfoot[C]{\small\thepage}
\renewcommand{\headrulewidth}{0.4pt}

% ── Section formatting ────────────────────────────────────────────────────────
\titleformat{\section}{\large\bfseries}{}{0em}{\thesection\quad}[\vspace{2pt}\titlerule]
\titleformat{\subsection}{\normalsize\bfseries}{}{0em}{\thesubsection\quad}
\titleformat{\subsubsection}{\normalsize\itshape\bfseries}{}{0em}{\thesubsubsection\quad}
\setcounter{secnumdepth}{3}

% ══════════════════════════════════════════════════════════════════════════════
\begin{document}
% ══════════════════════════════════════════════════════════════════════════════

% ── Title Page ────────────────────────────────────────────────────────────────
\begin{titlepage}
  \centering
  \vspace*{2cm}

  {\Huge\bfseries Wireless Drive-By-Wire\\[0.4em]
  \Large Design Document}

  \vspace{0.5cm}\rule{\linewidth}{0.5pt}\vspace{1cm}

  {\large Capstone Spring 2026}\\[0.4cm]
  {\normalsize ECE 4773 --- Senior Design}

  \vspace{2cm}
  \begin{tabular}{ll}
    \textbf{Document Version:} & 1.0 \\
    \textbf{Date:}             & May 2026 \\
    \textbf{Audience:}         & Future developers, maintainers \\
  \end{tabular}

  \vfill

  \begin{notebox}
    This document is intended for engineers who need to \textbf{repair, modify,
    or extend} the system.  It covers the internal hardware design, signal
    mappings, control algorithm, network protocol, and known design limitations.
    For operator instructions, see the separate \textit{User Manual}.
  \end{notebox}

  \vspace{1cm}
\end{titlepage}

\tableofcontents
\newpage

% ══════════════════════════════════════════════════════════════════════════════
\section{System Overview}
% ══════════════════════════════════════════════════════════════════════════════

The Wireless Drive-By-Wire (DBW) system replaces the mechanical inputs of a
Nissan Leaf (steering wheel, throttle, brake pedal) with electronically generated
analog voltages that the vehicle's ECU interprets as driver input.

An operator uses a Logitech G29 racing wheel and pedals connected to a laptop.
Control packets are transmitted over UDP through an Oracle Cloud relay server to an
ESP32 microcontroller mounted on a custom PCB inside the vehicle.  The ESP32 reads
the actual steering angle from the CAN bus and drives three systems:

\begin{itemize}
  \item \textbf{Steering} --- closed-loop PID, DAC voltages drive the power steering motor
  \item \textbf{Throttle} --- open-loop, DAC voltages increase from a baseline
  \item \textbf{Brake}    --- open-loop, DAC voltages decrease/increase from a baseline
\end{itemize}

\subsection{Top-Level Block Diagram}

\begin{center}
\renewcommand{\arraystretch}{1.4}
\small\ttfamily
\begin{tabular}{|c|c|c|c|c|}
  \hline
  \textrm{\textbf{G29}} & \textrm{USB} &
  \textrm{\textbf{Laptop}} & \textrm{UDP/internet} &
  \textrm{\textbf{OCI Relay Server}} \\
  \hline
  \multicolumn{5}{c}{$\downarrow$ WiFi UDP} \\
  \hline
  \multicolumn{2}{|c|}{\textrm{\textbf{ESP32 on PCB}}} &
  \multicolumn{1}{c|}{CAN bus $\leftarrow$} &
  \multicolumn{2}{c|}{\textrm{Nissan Leaf (CAN ID 0x2)}} \\
  \hline
  \multicolumn{2}{|c|}{DAC2 Ch C/D} &
  \multicolumn{1}{c|}{DAC1 Ch A/B} &
  \multicolumn{2}{c|}{DAC1 Ch C/D} \\
  \hline
  \multicolumn{2}{|c|}{Steering motor} &
  \multicolumn{1}{c|}{Throttle actuator} &
  \multicolumn{2}{c|}{Brake actuator} \\
  \hline
\end{tabular}
\end{center}

\subsection{File Map}

\begin{center}
\begin{tabular}{lll}
  \toprule
  \textbf{File} & \textbf{Runs on} & \textbf{Role} \\
  \midrule
  \texttt{g29\_to\_server\_original.py} & Operator laptop & Read G29, send UDP \\
  \texttt{relay\_server.py}             & OCI cloud VM    & Forward UDP packets \\
  \texttt{esp32\_car.ino}               & ESP32 on PCB    & Control actuators \\
  \bottomrule
\end{tabular}
\end{center}

% ══════════════════════════════════════════════════════════════════════════════
\section{Hardware Design}
% ══════════════════════════════════════════════════════════════════════════════

\subsection{Custom PCB}

The custom DBW PCB integrates the following onto a single board:

\begin{itemize}
  \item ESP32-WROOM-32 DevKit V1 (footprint / header)
  \item MCP2515 CAN bus controller (SPI) with 16\,MHz crystal
  \item Two MCP4728 quad 12-bit I$^2$C DACs
  \item Supporting passives (decoupling capacitors, pull-up resistors for I$^2$C)
\end{itemize}

The PCB was fabricated by JLCPCB (2-layer).

\subsection{ESP32 Pin Assignments}

\begin{center}
\begin{tabular}{llll}
  \toprule
  \textbf{ESP32 Pin} & \textbf{Signal} & \textbf{Interface} & \textbf{Connected to} \\
  \midrule
  18  & SPI SCK  & SPI  & MCP2515 SCK \\
  19  & SPI MISO & SPI  & MCP2515 SO \\
  23  & SPI MOSI & SPI  & MCP2515 SI \\
  5   & CAN CS   & SPI  & MCP2515 \(\overline{\text{CS}}\) \\
  4   & CAN INT  & GPIO & MCP2515 INT (active-low) \\
  21  & I2C1 SDA & I$^2$C bus 0 & DAC 1 (MCP4728 @ 0x60) \\
  22  & I2C1 SCL & I$^2$C bus 0 & DAC 1 (MCP4728 @ 0x60) \\
  33  & I2C2 SDA & I$^2$C bus 1 & DAC 2 (MCP4728 @ 0x60) \\
  32  & I2C2 SCL & I$^2$C bus 1 & DAC 2 (MCP4728 @ 0x60) \\
  \bottomrule
\end{tabular}
\end{center}

\begin{notebox}
Both DACs share the same I$^2$C address (\texttt{0x60}).  They are disambiguated
by being placed on \emph{separate} I$^2$C buses: \texttt{W1} (TwoWire instance 0)
and \texttt{W2} (TwoWire instance 1).  If you ever replace a DAC, verify that its
address pins (A0, A1, A2) are all grounded so it enumerates as \texttt{0x60}.
\end{notebox}

\subsection{MCP2515 CAN Controller}

The MCP2515 communicates with the ESP32 over SPI and with the vehicle over the
OBD-II CAN bus.

\begin{center}
\begin{tabular}{ll}
  \toprule
  \textbf{Parameter} & \textbf{Value} \\
  \midrule
  CAN bus speed    & 500\,kbps (\texttt{CAN\_500KBPS}) \\
  Crystal frequency & 16\,MHz (\texttt{MCP\_16MHZ}) \\
  Operating mode   & Normal (\texttt{MCP\_NORMAL}) \\
  Filter           & None (\texttt{MCP\_ANY} --- all IDs accepted) \\
  INT pin polarity & Active-low; triggers on any received frame \\
  \bottomrule
\end{tabular}
\end{center}

\begin{warningbox}
The firmware uses \texttt{MCP\_16MHZ}.  If the PCB crystal is ever replaced, it
\textbf{must} be 16\,MHz, not 8\,MHz.  A mismatch silently corrupts CAN timing;
the symptom is that \texttt{currentSTA} never updates and steering drives to its
limit.
\end{warningbox}

\subsection{MCP4728 DAC Configuration}

Both DACs are configured identically: internal voltage reference with $2\times$ gain.

\begin{center}
\begin{tabular}{lll}
  \toprule
  \textbf{Parameter} & \textbf{Setting} & \textbf{Effect} \\
  \midrule
  Voltage reference & \texttt{MCP4728\_VREF\_INTERNAL} & Uses 2.048\,V internal ref \\
  Gain              & \texttt{MCP4728\_GAIN\_2X}        & Full scale = $2\times2.048 = 4.096$\,V \\
  Resolution        & 12-bit                             & $4096$ steps, $\approx$1\,mV/step \\
  \bottomrule
\end{tabular}
\end{center}

These settings are written in \texttt{initDACChannels()} and saved to the DACs'
internal EEPROM via \texttt{saveToEEPROM()} so they survive power cycling.

\subsubsection*{Converting Voltage to DAC Count}

\[
  \text{DAC count} = \left\lfloor \frac{V_{\text{target}}}{4.096} \times 4095
  \right\rfloor
\]

The firmware stores the reference voltage as \texttt{float proper = 4.096} and
applies this formula in \texttt{updateBaseVoltages()} every loop iteration.

% ══════════════════════════════════════════════════════════════════════════════
\section{Signal Mapping and Baseline Voltages}
% ══════════════════════════════════════════════════════════════════════════════

\subsection{DAC Channel Assignments}

\begin{center}
\begin{tabular}{llllrr}
  \toprule
  \textbf{Variable} & \textbf{DAC} & \textbf{Ch} & \textbf{Function}
    & \textbf{Baseline V} & \textbf{DAC count} \\
  \midrule
  \texttt{steer1}   & DAC 2 & C & Steering main  & 2.47\,V & 2469 \\
  \texttt{steer2}   & DAC 2 & D & Steering sub   & 2.47\,V & 2469 \\
  \texttt{accel1Base} & DAC 1 & A & Throttle main & 0.37\,V & 369  \\
  \texttt{accel2Base} & DAC 1 & B & Throttle sub  & 0.75\,V & 749  \\
  \texttt{brake1Base} & DAC 1 & C & Brake main    & 3.42\,V & 3419 \\
  \texttt{brake2Base} & DAC 1 & D & Brake sub     & 1.51\,V & 1509 \\
  \bottomrule
\end{tabular}
\end{center}

\begin{notebox}
\texttt{updateBaseVoltages()} is called once per loop.  Although the values are
constant, keeping it in the loop makes them easy to change to live-tunable
variables in the future (e.g., read from a serial command).
\end{notebox}

\subsection{Steering Output (Push-Pull)}

The steering motor is driven push-pull: one DAC channel raises above midpoint while
the other drops by the same amount.  The motor controller interprets the voltage
\emph{difference} as torque magnitude and direction.

\begin{align*}
  \text{DAC 2, Ch C} &= \texttt{steer1} + \texttt{torque} \\
  \text{DAC 2, Ch D} &= \texttt{steer2} - \texttt{torque}
\end{align*}

At zero torque both channels output 2.47\,V and the motor is stationary.
\texttt{torque} is clamped to $\pm1500$ DAC counts ($\approx\pm1.5$\,V swing).

\subsection{Throttle Output}

\[
  \texttt{accelMag} = \lfloor\texttt{throttle} \times 300\rfloor \quad (0\text{--}300)
\]

\begin{align*}
  \text{DAC 1, Ch A} &= \texttt{accel1Base} + \texttt{accelMag} \\
  \text{DAC 1, Ch B} &= \texttt{accel2Base} + 2 \times \texttt{accelMag}
\end{align*}

Channel B receives twice the offset of channel A.  This mimics the dual-channel
pedal sensor redundancy scheme used by the Nissan Leaf's original pedal assembly
(primary and secondary sensors track each other with a 2:1 voltage ratio).

\subsection{Brake Output}

\[
  \texttt{brakeMag} = \lfloor\texttt{brake} \times 900\rfloor \quad (0\text{--}900)
\]

\begin{align*}
  \text{DAC 1, Ch C} &= \texttt{brake1Base} - \texttt{brakeMag} \\
  \text{DAC 1, Ch D} &= \texttt{brake2Base} + \texttt{brakeMag}
\end{align*}

\begin{warningbox}
The brake channels must \textbf{never} be written to 0\,V.  The vehicle ECU
monitors both brake sensor voltages; if either reads 0\,V it interprets this as
a broken wire (sensor fault) and may disable throttle or engage a failsafe.
The baseline voltages (3.42\,V and 1.51\,V) are the ``no braking'' idle state,
not 0\,V.
\end{warningbox}

\subsection{All Outputs Clamped}

Every DAC output is passed through \texttt{constrain(value, 0, 4095)} inside
\texttt{sendOutputs()} before being written.  This prevents integer overflow
from wrapping a count above 4095 back to a small positive value.

% ══════════════════════════════════════════════════════════════════════════════
\section{CAN Bus Interface}
% ══════════════════════════════════════════════════════════════════════════════

\subsection{Steering Angle Message}

The Nissan Leaf broadcasts the steering angle on CAN ID \texttt{0x2}.

\begin{center}
\begin{tabular}{llll}
  \toprule
  \textbf{Byte(s)} & \textbf{Type} & \textbf{Endian} & \textbf{Scaling} \\
  \midrule
  0--1 & \texttt{int16} & Little-endian & $\div 10$ = degrees \\
  2--7 & (unused) & --- & --- \\
  \bottomrule
\end{tabular}
\end{center}

Parsed in firmware as:

\begin{lstlisting}[style=cpp]
currentSTA = (int16_t)((buf[0]) | (buf[1] << 8)) / 10;
\end{lstlisting}

\texttt{currentSTA} is a signed 16-bit integer representing steering angle in
tenths of a degree, divided to give degrees directly.  Positive values are
right-of-center; negative values are left-of-center (confirm the sign convention
during initial vehicle testing).

\subsection{Reading Strategy}

The firmware polls the MCP2515 INT pin (active-low, GPIO 4) every loop.
When the pin is low, a frame is waiting; the firmware reads and filters for
ID \texttt{0x2}.  All other IDs are silently discarded (\texttt{MCP\_ANY} mode
accepts everything, but only \texttt{0x2} is stored).

\begin{notebox}
If you need to capture additional CAN signals (speed, gear position, etc.),
add more \texttt{if (id == 0xXX)} branches inside \texttt{readCAN()}.  No
filter reconfiguration is needed because the firmware already accepts all IDs.
\end{notebox}

% ══════════════════════════════════════════════════════════════════════════════
\section{Steering PID Control Loop}
% ══════════════════════════════════════════════════════════════════════════════

\subsection{Loop Structure}

The control loop runs in the Arduino \texttt{loop()} function at the MCU's
natural speed (no fixed rate delay).  Loop timing is measured via
\texttt{micros()} and passed as \texttt{dt} (seconds) into the PID.  Values of
\texttt{dt} outside $(0, 1)$ are clamped to 0.01\,s to prevent divide-by-zero
or derivative spikes on startup.

\subsection{Target Calculation and Smoothing}

\[
  \texttt{targetSTA} = \texttt{steering} \times 250 \quad (\text{degrees})
\]

A first-order IIR low-pass filter smooths the target before it enters the PID:

\[
  \texttt{smoothTarget}_{k} = \texttt{smoothTarget}_{k-1}
    + \alpha \cdot (\texttt{targetSTA} - \texttt{smoothTarget}_{k-1})
\]

where $\alpha = 10 \cdot dt$.  At a typical loop rate of $\sim$1\,ms,
$\alpha \approx 0.01$, giving a time constant of $\tau = 1/10 = 100$\,ms.

\begin{notebox}
The smoothing filter prevents abrupt steering commands (e.g., from a dropped and
re-received packet with a large position jump) from immediately spiking the PID
output.  Reducing the coefficient (e.g., $5 \cdot dt$) makes steering feel
``heavier''; increasing it makes it more responsive.
\end{notebox}

\subsection{PID Gains}

\begin{center}
\begin{tabular}{llll}
  \toprule
  \textbf{Gain} & \textbf{Variable} & \textbf{Value} & \textbf{Units} \\
  \midrule
  Proportional & \texttt{kp} & 3.6  & DAC counts / degree \\
  Integral     & \texttt{ki} & 0.05 & DAC counts / (degree$\cdot$s) \\
  Derivative   & \texttt{kd} & 0.1  & DAC counts$\cdot$s / degree \\
  \bottomrule
\end{tabular}
\end{center}

\subsection{Anti-Windup and Clamping}

The firmware implements two windup-prevention measures:

\begin{enumerate}
  \item \textbf{Dead band reset}: if $|\text{error}| < 2°$, the integral is reset
        to zero.  This prevents integrator drift from accumulated small errors near
        the target.
  \item \textbf{Conditional integration}: the integral accumulates only when the
        proportional term alone has not saturated the output (i.e.,
        $|k_p e + k_i I| < \texttt{maxTorque}$).
  \item \textbf{Integral clamp}: the integral is hard-clamped to $\pm250$
        DAC counts after each update.
\end{enumerate}

The final output is clamped to $\pm1500$ DAC counts (\texttt{maxTorque}):

\[
  \texttt{torque} = \text{clamp}\!\left(
    k_p e + k_i I + k_d \dot{e},\;
    -1500,\; +1500
  \right)
\]

\subsection{PID Tuning Guide}

If you modify the vehicle or steering motor, you will likely need to re-tune.
Follow this sequence on a \textbf{stationary vehicle}:

\begin{enumerate}
  \item Set \texttt{ki = 0} and \texttt{kd = 0}.  Increase \texttt{kp} until
        steering reaches the target but oscillates.
  \item Back off \texttt{kp} by $\sim$20\%.
  \item Increase \texttt{kd} until oscillation is damped.
  \item Restore a small \texttt{ki} (start at 0.01) to eliminate steady-state error.
  \item Validate at low speed before testing at higher speeds.
\end{enumerate}

% ══════════════════════════════════════════════════════════════════════════════
\section{Network Architecture}
% ══════════════════════════════════════════════════════════════════════════════

\subsection{Why a Relay Server?}

The ESP32 connects to a mobile hotspot.  Mobile carriers perform Network Address
Translation (NAT): the ESP32 has a private IP on the hotspot LAN and a shared
public IP assigned to the carrier's gateway.  It is unreachable from the internet
by direct address.

The relay server solves this via \textbf{UDP hole-punching}:

\begin{enumerate}
  \item The ESP32 sends a \texttt{HEARTBEAT} packet to the OCI server every 2\,s.
        Because the packet originates from the ESP32, the hotspot's NAT table
        creates a mapping: \textit{(public IP : ephemeral port)} $\leftrightarrow$
        \textit{ESP32 private IP}.
  \item The relay server records the ESP32's source address and port from the
        arriving \texttt{HEARTBEAT}.
  \item When a \texttt{CMD} packet arrives from the operator laptop, the relay
        forwards it to that recorded address.  The hotspot's NAT entry still
        exists (kept alive by the 2\,s heartbeat), so the packet is delivered.
\end{enumerate}

\subsection{Relay Server Logic}

\begin{lstlisting}[style=python]
if message == "HEARTBEAT":
    car_address = addr          # register / refresh ESP32 endpoint
elif car_address and addr != car_address:
    sock.sendto(data, car_address)  # forward operator packet to car
\end{lstlisting}

Any packet originating \emph{from} the car (\texttt{addr == car\_address}) that
is not a \texttt{HEARTBEAT} is silently dropped.  This prevents the car from
accidentally feeding data back to itself if it ever sends non-heartbeat traffic.

\subsection{Packet Protocol}

All packets are plain UTF-8 text, terminated with the UDP packet boundary
(no newline needed).

\begin{center}
\begin{tabular}{lll}
  \toprule
  \textbf{Direction} & \textbf{Format} & \textbf{Notes} \\
  \midrule
  ESP32 $\to$ Server & \texttt{HEARTBEAT} & Sent every 2\,s \\
  Laptop $\to$ Server & \texttt{CMD;\textit{seq};\textit{steer};\textit{throttle};\textit{brake}}
    & Sent at 60\,Hz \\
  Server $\to$ ESP32 & \textit{(forwarded CMD verbatim)} & No modification \\
  \bottomrule
\end{tabular}
\end{center}

\begin{center}
\begin{tabular}{llll}
  \toprule
  \textbf{Field} & \textbf{Type} & \textbf{Range} & \textbf{Description} \\
  \midrule
  \texttt{seq}      & uint32  & 0, 1, 2\ldots & Monotonic counter; starts at 0 \\
  \texttt{steer}    & float32 & $-1.0\ldots+1.0$ & $-1.0=$ full left \\
  \texttt{throttle} & float32 & $0.0\ldots1.0$   & $0.0=$ released \\
  \texttt{brake}    & float32 & $0.0\ldots1.0$   & $0.0=$ released \\
  \bottomrule
\end{tabular}
\end{center}

The ESP32 parses with \texttt{sscanf}:

\begin{lstlisting}[style=cpp]
sscanf(packetBuffer, "CMD;%lu;%f;%f;%f", &seq, &s, &t, &b)
\end{lstlisting}

There is no sequence-number validation or out-of-order rejection in the current
firmware.  If a late packet arrives it will overwrite the current control state.

\subsection{Network Configuration}

\begin{center}
\begin{tabular}{ll}
  \toprule
  \textbf{Parameter} & \textbf{Value} \\
  \midrule
  Relay server public IP & \texttt{147.224.143.221} \\
  UDP port (all parties) & \texttt{4210} \\
  Sender rate            & 60\,Hz \\
  Heartbeat interval     & 2\,s \\
  \bottomrule
\end{tabular}
\end{center}

% ══════════════════════════════════════════════════════════════════════════════
\section{Software Architecture}
% ══════════════════════════════════════════════════════════════════════════════

\subsection{\texttt{g29\_to\_server\_original.py} (Operator Laptop)}

\begin{center}
\begin{tabular}{ll}
  \toprule
  \textbf{Constant}          & \textbf{Value / Description} \\
  \midrule
  \texttt{STEER\_AXIS\_INDEX}    & 0 --- steering wheel axis \\
  \texttt{THROTTLE\_AXIS\_INDEX} & 2 --- throttle pedal axis \\
  \texttt{BRAKE\_AXIS\_INDEX}    & 3 --- brake pedal axis \\
  \texttt{SERVER\_IP}           & \texttt{147.224.143.221} \\
  \texttt{SERVER\_PORT}         & 4210 \\
  \texttt{REFRESH\_HZ}          & 60 \\
  \bottomrule
\end{tabular}
\end{center}

The G29 pedal axes report \texttt{+1.0} when \emph{released} and \texttt{-1.0}
when \emph{fully pressed}.  The sender normalizes and inverts:

\begin{lstlisting}[style=python]
def normalize_01_from_minus1_1(val: float) -> float:
    return max(0.0, min(1.0, (val + 1.0) / 2.0))

throttle = 1.0 - normalize_01_from_minus1_1(joystick.get_axis(2))
brake    = 1.0 - normalize_01_from_minus1_1(joystick.get_axis(3))
\end{lstlisting}

Step-by-step for the \textbf{throttle pedal fully pressed}
(axis reports \texttt{-1.0}):
\begin{enumerate}
  \item \texttt{normalize\_01(-1.0)} $= ((-1)+1)/2 = 0.0$
  \item \texttt{throttle} $= 1.0 - 0.0 = 1.0$
\end{enumerate}

Step-by-step for the \textbf{throttle pedal released}
(axis reports \texttt{+1.0}):
\begin{enumerate}
  \item \texttt{normalize\_01(+1.0)} $= ((+1)+1)/2 = 1.0$
  \item \texttt{throttle} $= 1.0 - 1.0 = 0.0$
\end{enumerate}

Steering is passed through \emph{without} normalization; the G29 steering axis
already reports \texttt{-1.0} (full left) to \texttt{+1.0} (full right).

\subsection{\texttt{relay\_server.py} (OCI Cloud VM)}

Minimal single-socket server.  Runs on \texttt{0.0.0.0:4210}.  The only state
it maintains is \texttt{car\_address} --- the most recently seen
\texttt{HEARTBEAT} source.  All non-HEARTBEAT packets arriving from an address
other than \texttt{car\_address} are forwarded verbatim.

\begin{designbox}
The relay server does not time out \texttt{car\_address}.  If the ESP32
reboots and obtains a new public IP/port from the hotspot, the first new
\texttt{HEARTBEAT} will update the stored address.  There is no danger of
sending packets to a stale address indefinitely.
\end{designbox}

The server also prints a \texttt{DEBUG} line for every received packet.
This is intentional for initial development but will flood the console during
normal operation.  To silence it, remove or comment out the two
\texttt{print(f"DEBUG...")} lines.

\subsection{\texttt{esp32\_car.ino} (ESP32 on PCB)}

\subsubsection*{Setup Sequence}

\begin{enumerate}
  \item \texttt{Serial.begin(115200)}
  \item Initialize I$^2$C buses (W1 on 21/22, W2 on 33/32, both at 400\,kHz)
  \item Initialize DAC 1 and DAC 2 (\texttt{begin(0x60, \&W1/W2)})
  \item \texttt{initDACChannels()} --- write initial DAC values + save to EEPROM
  \item \texttt{initCAN()} --- start SPI and MCP2515 at 500\,kbps / 16\,MHz
  \item \texttt{initWiFi()} --- connect to hotspot, start UDP listener on port 4210
\end{enumerate}

\subsubsection*{Main Loop Sequence (per iteration)}

\begin{enumerate}
  \item Compute \texttt{dt} from \texttt{micros()} delta
  \item Send \texttt{HEARTBEAT} if $>2000$\,ms since last heartbeat
  \item \texttt{receiveUDP()} --- parse latest \texttt{CMD} packet into
        \texttt{steering}, \texttt{throttle}, \texttt{brake}
  \item \texttt{readCAN()} --- update \texttt{currentSTA} from CAN ID \texttt{0x2}
  \item \texttt{updateBaseVoltages()} --- recompute DAC baseline counts
  \item Constrain \texttt{steering}, \texttt{throttle}, \texttt{brake}
  \item Compute \texttt{targetSTA}, apply low-pass filter to \texttt{smoothTarget}
  \item Compute \texttt{accelMag} and \texttt{brakeMag}
  \item \texttt{calcTorque(smoothTarget, currentSTA, dt)} --- run PID
  \item \texttt{sendOutputs()} --- write all six DAC values
  \item Print debug to Serial every 100\,ms
\end{enumerate}

% ══════════════════════════════════════════════════════════════════════════════
\section{Known Limitations and Design Decisions}
% ══════════════════════════════════════════════════════════════════════════════

\subsection{No Watchdog / Failsafe}

\begin{warningbox}
There is \textbf{no watchdog timer} in the current firmware.  If the operator
laptop stops sending packets (network drop, crash, Ctrl+C), the ESP32 continues
targeting the \textbf{last received} steering angle, throttle, and brake values
indefinitely.  This means the car will continue accelerating and steering at its
last commanded state until the safety driver intervenes.
\end{warningbox}

A simple fix is to timestamp the last received packet and reset all control
variables to zero if no packet arrives within a timeout:

\begin{lstlisting}[style=cpp]
#define WATCHDOG_MS 200
unsigned long lastPacketMs = 0;

// In receiveUDP(), after successful parse:
lastPacketMs = millis();

// In loop(), after receiveUDP():
if (millis() - lastPacketMs > WATCHDOG_MS) {
    steering = 0.0f;
    throttle = 0.0f;
    brake    = 0.0f;
    integral = 0.0f;  // reset PID integrator
}
\end{lstlisting}

\subsection{Open-Loop Throttle and Brake}

Throttle and brake are driven open-loop: there is no speed sensor feedback.
The DAC offsets were determined empirically on the test vehicle.  If the vehicle
or actuators change, the scaling factors (300 for throttle, 900 for brake) will
need to be re-calibrated.

\subsection{EEPROM Initialization Mismatch}

\texttt{initDACChannels()} writes initial values to EEPROM that do \emph{not}
match the runtime channel assignments.  For example, it writes steering midpoint
values to DAC\,1 channels C/D, but during runtime DAC\,1 Ch\,C/D drive the
brake.  This is harmless because the firmware immediately overwrites all channels
via \texttt{fastWrite()} on the first loop iteration.  However, if the ESP32
powers up without running \texttt{loop()} (e.g., it crashes in \texttt{setup()}),
the EEPROM-default outputs may be unexpected.  A future revision should align the
EEPROM values with the runtime assignments.

\subsection{DEBUG Prints in relay\_server.py}

The relay server prints a debug line for \emph{every} received UDP packet:

\begin{lstlisting}[style=terminal]
DEBUG - Received: 'CMD;0;0.0000;0.0000;0.0000' from ('1.2.3.4', 54321)
\end{lstlisting}

At 60\,Hz this generates 60 lines per second and will fill terminal scrollback
quickly.  Remove the two \texttt{print(f"DEBUG...")} lines once the system is
confirmed working.

\subsection{Relay Server Has No Timeout on Car Address}

If the car physically moves out of hotspot range, \texttt{car\_address} in the
relay server is never cleared.  CMD packets from the operator will be silently
forwarded to a stale address with no error shown.  The operator only notices
because the car stops responding.  A future revision should print a warning
if no \texttt{HEARTBEAT} is received within $\sim$10\,s.

% ══════════════════════════════════════════════════════════════════════════════
\section{Modification Guide}
% ══════════════════════════════════════════════════════════════════════════════

\subsection{Changing WiFi Credentials}

In \texttt{esp32\_car.ino}, update and re-flash:

\begin{lstlisting}[style=cpp]
const char* ssid     = "NEW_SSID";
const char* password = "NEW_PASSWORD";
\end{lstlisting}

\subsection{Changing the Relay Server IP or Port}

\textbf{In \texttt{esp32\_car.ino}:}
\begin{lstlisting}[style=cpp]
const char* serverIP = "NEW.IP.ADDRESS";
const int   serverPort = NEW_PORT;
\end{lstlisting}

\textbf{In \texttt{g29\_to\_server\_original.py}:}
\begin{lstlisting}[style=python]
SERVER_IP   = "NEW.IP.ADDRESS"
SERVER_PORT = NEW_PORT
\end{lstlisting}

Both files must be updated.  The relay server binds to \texttt{0.0.0.0} and uses
whatever port is defined in \texttt{UDP\_PORT}, so update that constant too if
you change the port.

\subsection{Adjusting Throttle and Brake Scaling}

The scaling factors are literal constants in \texttt{loop()}:

\begin{lstlisting}[style=cpp]
accelMag = pedalMagnitude(appliedThrottle, 300.0f);  // max offset = 300
brakeMag = pedalMagnitude(appliedBrake,    900.0f);  // max offset = 900
\end{lstlisting}

Increase the second argument to increase the maximum actuator travel.
Verify that the resulting DAC count (\texttt{accel1Base + accelMag}) does not
exceed 4095 at full input, or the \texttt{constrain()} in \texttt{sendOutputs()}
will clip it.

\subsection{Reading Additional CAN Signals}

Add cases inside \texttt{readCAN()}:

\begin{lstlisting}[style=cpp]
if (id == 0x2) {
    currentSTA = (int16_t)((buf[0]) | (buf[1] << 8)) / 10;
}
// Example: read vehicle speed from a different CAN ID
if (id == 0x1B4) {
    vehicleSpeed = (buf[0] | (buf[1] << 8)) * 0.01f;  // example scaling
}
\end{lstlisting}

No filter configuration is needed because the MCP2515 is in \texttt{MCP\_ANY}
mode.

\subsection{Adding the Watchdog}

See Section~9.1 for the recommended code snippet.  After adding it, test by
unplugging the laptop's network cable with the car stationary and confirming
that all actuators return to their idle state within 200\,ms.

\end{document}
